%======================================================================

\subsection{Geodetic Theory and Coordinate Systems}
\label{sec:geodetic_theory}

The following subsections provide the theoretical foundations for the coordinate transformations used throughout this analysis. The goal is to make the paper self-contained for a physicist who may not be familiar with the conventions and subtleties of geodetic reference frames.

\subsubsection{The WGS\,84 Ellipsoid}
\label{sec:wgs84}

The World Geodetic System 1984 (WGS\,84) defines a geocentric reference ellipsoid---an oblate spheroid of revolution---as the fundamental surface for GPS positioning~\cite{nga2014}. The defining parameters are the semi-major axis $a$ and the flattening $f$:
\begin{align}
a &= 6\,378\,137.0~\text{m}, \label{eq:wgs84_a}\\
f &= \frac{1}{298.257223563}. \label{eq:wgs84_f}
\end{align}
From these, the semi-minor axis $b = a(1-f) = 6\,356\,752.314~\text{m}$ and the first eccentricity squared $e^2 = 2f - f^2 = 0.00669437999014$ follow immediately.

A point near Earth's surface is described in geodetic coordinates $(\varphi, \lambda, h)$, where $\varphi$ is the geodetic latitude (the angle between the ellipsoid normal and the equatorial plane), $\lambda$ is the longitude (measured east from the Greenwich meridian), and $h$ is the ellipsoidal height (the distance above the ellipsoid along its outward normal). The ellipsoidal height is \emph{not} the elevation above mean sea level; the distinction is addressed in Section~\ref{sec:geoid}.

\paragraph{Geodetic to ECEF Cartesian.} The transformation from geodetic coordinates $(\varphi, \lambda, h)$ to geocentric Earth-centered Earth-fixed (ECEF) Cartesian coordinates $(X, Y, Z)$ is~\cite{nga2014}:
\begin{align}
X &= (N + h)\cos\varphi\cos\lambda, \label{eq:geod2xyz_X}\\
Y &= (N + h)\cos\varphi\sin\lambda, \label{eq:geod2xyz_Y}\\
Z &= \bigl[N(1 - e^2) + h\bigr]\sin\varphi, \label{eq:geod2xyz_Z}
\end{align}
where $N$ is the radius of curvature in the prime vertical:
\begin{equation}
N(\varphi) = \frac{a}{\sqrt{1 - e^2 \sin^2\!\varphi}}.
\label{eq:radius_N}
\end{equation}
At the TA site ($\varphi \approx 39.3\degree$), $N \approx 6{,}381{,}040$\,m.

\paragraph{ECEF Cartesian to geodetic.} The inverse transformation---given $(X, Y, Z)$, find $(\varphi, \lambda, h)$---has an exact solution for $\lambda$:
\begin{equation}
\lambda = \arctan\!\left(\frac{Y}{X}\right),
\label{eq:xyz2geod_lam}
\end{equation}
but $\varphi$ and $h$ are coupled through Eq.~\eqref{eq:radius_N} and must be solved iteratively. Bowring's method~\cite{bowring1976} provides rapid convergence, typically reaching sub-millimeter accuracy in a single iteration. The algorithm proceeds as follows. Define $p = \sqrt{X^2 + Y^2}$. Compute an initial estimate of the reduced latitude:
\begin{equation}
\beta_0 = \arctan\!\left(\frac{Z}{p\,(1 - e^2)}\right).
\label{eq:bowring_beta0}
\end{equation}
Then iterate (usually once):
\begin{align}
\varphi &= \arctan\!\left(\frac{Z + e'^2\,b\,\sin^3\!\beta}{p - e^2\,a\,\cos^3\!\beta}\right), \label{eq:bowring_phi}\\
\beta &= \arctan\!\left(\frac{(1-f)\sin\varphi}{\cos\varphi}\right), \label{eq:bowring_beta}
\end{align}
where $e'^2 = e^2/(1-e^2)$ is the second eccentricity squared. After convergence, the ellipsoidal height is:
\begin{equation}
h = \frac{p}{\cos\varphi} - N(\varphi).
\label{eq:bowring_h}
\end{equation}
This is the method used in the survey processing code (\texttt{survey\_processor.c}) and independently verified via the Python script \texttt{xyz2csv.py}. The self-consistency check in Table~\ref{tab:xyz_validation} confirms agreement to within 3\,mm.

\subsubsection{The Geoid and Orthometric Height}
\label{sec:geoid}

The \emph{geoid} is the equipotential surface of Earth's gravity field that best approximates mean sea level in the open ocean~\cite{torge2012}. Because the mass distribution inside the Earth is not uniform, the geoid is an irregular surface that undulates above and below the reference ellipsoid by up to $\pm 100$\,m globally. The geoid undulation $\mathcal{N}$ (we use calligraphic $\mathcal{N}$ to distinguish it from the radius of curvature $N$) at a point is the signed distance from the ellipsoid to the geoid along the ellipsoid normal.

The fundamental height relationship is:
\begin{equation}
h = H + \mathcal{N},
\label{eq:height_relation}
\end{equation}
where $h$ is the ellipsoidal height (from GPS), $H$ is the orthometric height (physical elevation above the geoid, i.e., above ``sea level''), and $\mathcal{N}$ is the geoid undulation. Rearranging: $H = h - \mathcal{N}$.

In the United States, orthometric heights are referenced to the North American Vertical Datum of 1988 (NAVD\,88), and the geoid model GEOID12B (or its successors) provides the undulation $\mathcal{N}$ on a grid~\cite{roman2012}. At the Telescope Array site, $\mathcal{N} \approx -19.9$\,m. The negative sign means the geoid lies \emph{below} the WGS\,84 ellipsoid at this location, so that orthometric heights (elevations) are approximately 19.9\,m higher than ellipsoidal heights. For the BR1 monument, as an example: $h_{\text{WGS84}} = 1395.053$\,m, $\mathcal{N} \approx -19.9$\,m, and $H_{\text{ortho}} \approx 1395.053 - (-19.9) = 1415.0$\,m, consistent with the OPUS-reported NAVD\,88 orthometric height of 1415.778\,m (the small residual reflects the geoid model precision).

It is critical to distinguish the geoid undulation from the NAD\,83/ITRF height difference. The $\sim\!0.73$\,m offset documented in this paper is a difference between two ellipsoidal heights in two different geodetic datums; it has nothing to do with the $\sim\!19.9$\,m geoid undulation, which converts between ellipsoidal and orthometric height within a single datum.

\subsubsection{Deflection of the Vertical}
\label{sec:deflection}

At any point on the Earth's surface, the direction of the local gravity vector (the plumb line) generally does not coincide with the normal to the reference ellipsoid. The angular difference is called the \emph{deflection of the vertical}~\cite{torge2012}. It is decomposed into two components:
\begin{itemize}[nosep]
\item $\xi$ (xi): the north--south component, defined as positive when the gravity vector is deflected toward the north relative to the ellipsoid normal.
\item $\eta$ (eta): the east--west component, defined as positive when the gravity vector is deflected toward the east.
\end{itemize}
Typical magnitudes are a few arcseconds (1\,arcsec $\approx$ 4.85\,$\mu$rad $\approx$ 31\,mm/km). At the TA sites, the deflection parameters range from $-4.3''$ to $+1.8''$ (Table~\ref{tab:deflection}).

\paragraph{Why this matters for total station observations.} A total station instrument levels itself using a compensator that aligns to the local gravity vector---the plumb line. All angles measured by the instrument (azimuths and elevations) are therefore referenced to the \emph{astronomic} (gravity-aligned) local frame, not the \emph{geodetic} (ellipsoid-normal) frame. The difference between these two frames is precisely the deflection of the vertical. If uncorrected, the deflection introduces systematic errors in the computed target positions that grow linearly with distance: at a range of 100\,m, a 4\,arcsec deflection produces a lateral error of approximately 2\,mm, and at 700\,m (the MD1--MD2 baseline), approximately 14\,mm.

\paragraph{The deflection rotation.} To convert observations from the astronomic local frame to the geodetic local frame, a small rotation is applied. Let $(\delta E_a, \delta N_a, \delta U_a)$ be the East--North--Up offsets in the astronomic frame (as measured by the total station). The geodetic offsets $(\delta E_g, \delta N_g, \delta U_g)$ are obtained by applying the deflection-of-vertical rotation~\cite{torge2012}:
\begin{equation}
\begin{pmatrix} \delta E_g \\ \delta N_g \\ \delta U_g \end{pmatrix}
=
\begin{pmatrix}
1 & 0 & \eta \\
0 & 1 & -\xi \\
-\eta & \xi & 1
\end{pmatrix}
\begin{pmatrix} \delta E_a \\ \delta N_a \\ \delta U_a \end{pmatrix},
\label{eq:deflection_rotation}
\end{equation}
where $\xi$ and $\eta$ are in radians. This is a first-order (small-angle) rotation; the full rotation matrix differs only at second order in $(\xi, \eta)$, which is negligible for arcsecond-level deflections. The correction is applied in the survey processing code (Section~\ref{sec:survey_math}) after converting each total station observation to local ENU offsets and before transforming to ECEF.

\subsubsection{The NAD\,83 to ITRF Transformation}
\label{sec:helmert}

The transformation between the NAD\,83 and ITRF reference frames is a 14-parameter Helmert (similarity) transformation maintained by NGS~\cite{soler2004}. The general form, for a position vector $\mathbf{X}$ at epoch $t$ relative to a reference epoch $t_0$, is:
\begin{equation}
\mathbf{X}_{\text{ITRF}}(t) = \mathbf{X}_{\text{NAD83}} + \mathbf{T}(t_0) + \delta s(t_0)\,\mathbf{X}_{\text{NAD83}} + \mathbf{R}(t_0)\,\mathbf{X}_{\text{NAD83}} + \dot{\mathbf{T}}\,(t - t_0) + \delta\dot{s}\,(t - t_0)\,\mathbf{X}_{\text{NAD83}} + \dot{\mathbf{R}}\,(t - t_0)\,\mathbf{X}_{\text{NAD83}},
\label{eq:helmert}
\end{equation}
where:
\begin{itemize}[nosep]
\item $\mathbf{T} = (T_x, T_y, T_z)^{\!\top}$ is the translation vector (three parameters),
\item $\delta s$ is a differential scale factor (one parameter),
\item $\mathbf{R}$ is an antisymmetric rotation matrix encoding three small rotation angles $(\omega_x, \omega_y, \omega_z)$,
\item $\dot{\mathbf{T}}$, $\delta\dot{s}$, $\dot{\mathbf{R}}$ are the time derivatives of the above (seven rate parameters).
\end{itemize}
The 14 parameters ($\mathbf{T}$, $\delta s$, $\mathbf{R}$ and their rates) are published by NGS for each ITRF realization. For ITRF2000 relative to NAD\,83(CORS96), the translations are of order 1\,m, the rotations are of order 1\,mas (milliarcsecond), and the scale difference is of order $10^{-9}$.

At the Telescope Array site, the net effect of this transformation on the height coordinate reduces to the $\sim\!0.73$\,m offset documented in Eq.~\eqref{eq:frame_offset}. The horizontal offsets are a few centimeters. The precise value varies weakly with latitude (a few mm over the $\sim\!35$\,km north--south extent of the FD triangle), which accounts for the marker-to-marker variation in $\Delta h$ seen in Table~\ref{tab:height_corrections}.

\subsubsection{OPUS Static Positioning}
\label{sec:opus}

The Online Positioning User Service (OPUS) is a free service operated by the U.S.\ National Geodetic Survey (NGS)~\cite{soler2003}. The user submits a RINEX file containing dual-frequency GPS observations from a static occupation, and OPUS returns precise coordinates in both NAD\,83 and ITRF.

The processing strategy is carrier-phase double-differencing. OPUS selects three nearby Continuously Operating Reference Stations (CORS) and forms double-differenced observations, which cancel satellite and receiver clock errors and greatly reduce atmospheric effects. The key processing steps are:
\begin{enumerate}[nosep]
\item \textbf{CORS selection:} Three reference stations are selected based on proximity, data quality, and geometric strength (low PDOP). Baselines to the user station are typically 50--300\,km.
\item \textbf{Double-differencing:} For each epoch, the carrier-phase observable is differenced between satellites (eliminating receiver clock) and between stations (eliminating satellite clock). The resulting double-differenced observable is sensitive only to the baseline geometry, integer ambiguities, and residual atmospheric delays.
\item \textbf{Ambiguity resolution:} The carrier-phase observations are ambiguous by an integer number of wavelengths ($\lambda_{L1} = 19.0$\,cm, $\lambda_{L2} = 24.4$\,cm). OPUS resolves these integer ambiguities using a combination of wide-lane and narrow-lane techniques. The percentage of resolved ambiguities is reported as a quality metric; values below 80\% indicate a degraded solution.
\item \textbf{Position estimation:} With ambiguities fixed, the baseline vector to each CORS is determined by least-squares adjustment. The user position is the average of the three independently determined baseline solutions.
\item \textbf{Coordinate output:} The final position is reported in ITRF at the mean epoch of observation, and in NAD\,83(CORS96)(Epoch:2002.0000) after applying the Helmert transformation of Section~\ref{sec:helmert}.
\end{enumerate}
For a 6-hour or longer static occupation with a survey-grade antenna, OPUS typically achieves horizontal precision of 1--2\,cm and vertical precision of 2--5\,cm (95\% confidence) per session. The inverse-variance weighted averages of multiple sessions used in this work (Table~\ref{tab:height_corrections}) reduce the vertical uncertainty to 4--20\,mm.


%======================================================================
